\documentclass[a4paper, 12pt]{article}



%%% Работа с русским языком
\usepackage{cmap}					% поиск в PDF
\usepackage{mathtext} 				% русские буквы в формулах
\usepackage[T2A]{fontenc}			% кодировка
\usepackage[utf8]{inputenc}			% кодировка исходного текста
\usepackage[russian]{babel}	% локализация и переносы

%%% Дополнительная работа с математикой
\usepackage{amsmath,amsfonts,amssymb,amsthm,mathtools} % AMS
\usepackage{icomma} % "Умная" запятая: $0,2$ --- число, $0, 2$ --- перечисление

%% Номера формул
%\mathtoolsset{showonlyrefs=true} % Показывать номера только у тех формул, на которые есть \eqref{} в тексте.

%% Шрифты
\usepackage{euscript}	 % Шрифт Евклид
\usepackage{mathrsfs} % Красивый матшрифт

%% Поля
\usepackage[left=2cm,right=2cm,top=2cm,bottom=2cm,bindingoffset=0cm]{geometry}

%% Русские списки
\usepackage{enumitem}
\makeatletter
\AddEnumerateCounter{\asbuk}{\russian@alph}{щ}
\makeatother

%%% Работа с картинками
\usepackage{graphicx}  % Для вставки рисунков
\graphicspath{{images/}{images2/}}  % папки с картинками
\setlength\fboxsep{3pt} % Отступ рамки \fbox{} от рисунка
\setlength\fboxrule{1pt} % Толщина линий рамки \fbox{}
\usepackage{wrapfig} % Обтекание рисунков и таблиц текстом

%%% Работа с таблицами
\usepackage{array,tabularx,tabulary,booktabs} % Дополнительная работа с таблицами
\usepackage{longtable}  % Длинные таблицы
\usepackage{multirow} % Слияние строк в таблице
\usepackage[table,xcdraw]{xcolor} % Цветные таблицы

%% Красная строка
\setlength{\parindent}{2em}

%% Интервалы
\linespread{1}
\usepackage{multirow}

%% TikZ
\usepackage{tikz}
\usetikzlibrary{graphs,graphs.standard}

%% Верхний колонтитул
\usepackage{fancyhdr}
\pagestyle{fancy}

%% Перенос знаков в формулах (по Львовскому)
\newcommand*{\hm}[1]{#1\nobreak\discretionary{}
	{\hbox{$\mathsurround=0pt #1$}}{}}

%% Мои дополнения
\usepackage{float} %Добавляет возможность работы с командой [H] которая улучшает расположение на странице
\usepackage{gensymb} %Красивые градусы
\usepackage{graphicx}               % Импорт изображений
\usepackage{caption} % Пакет для подписей к рисункам, в частности, для работы caption*
\usepackage{indentfirst} % для отступа красной строки после заголовка


\begin{document}
\newcommand{\HRule}{\rule{\linewidth}{0.7mm}} % Defines a new command for the horizontal lines, change thickness here
	
	\begin{center}
		\large\textbf{Московский Физико-Технический Институт}\\
		\large\textbf{(государственный университет)}
	
		\vfill
		
		\Large Тема проекта <<телеграмм-бот для физической культуры>>\\[0.5cm] % Minor heading such as course title
		
		%----------------------------------------------------------------------------------------
		%	TITLE SECTION
		%----------------------------------------------------------------------------------------
		
		\HRule
		\\[0.4cm]
		{ \huge \bfseries Компьютерные технологии}
		\\[0.4cm] % Title of your document
		\HRule
		\\[0.5cm]
		
		\ \\
	\textbf{\large Автор:} \\	
	\large Капылов Максим Б01-001

		\vfill
		\hspace*{-0.8 cm}\includegraphics[width=100 pt]{frkt_logo}\\
		\large Долгопрудный, 2022
	\end{center}

\newpage
\setcounter{page}{2}
\fancyfoot[c]{\thepage}
\fancyhead[L] {Компьютерные технологии}
\fancyhead[R]{}
\newpage 

\tableofcontents

\newpage
\section{Введение}
 Идея проекта заключается в том, чтобы написать бота, который бы упрощал работу с учащимися в телеграмме. В этом проекте будет использован \textbf{python}, модуль \textbf{aiogram} и другое. Будет использована база данных (\textbf{sqlite3}) для хранения информации о пользователях и админах.  Бот будет разделен на 3 части:
 \begin{enumerate}
 \item[•] Админская чать
 \item[•] Клиентская чать
 \item[•] Работа с другими ботами по управлению телеграмм-каналами
 \end{enumerate}
 \section{Действия перед началом}
 Стоит это идею обсудить с преподавателями физической культуры. Также проблема в том, что у каждого преподавателя есть свой телеграмм-канал и им будет неудобно совместно управлять каналом.
 
  Составление плана интерфейса для учащихся и администраторов. Изучение функционала самого телеграма (я им почти не пользуюсь). Также \textbf{найти других участников для проекта}.
  \section{Основной функционал}
 Самое главное, что бот должен уметь - это создавать посты, под которыми будут отмечатся учащиеся. Текст поста примерно такой:  \textbf{17.05 пятница 18.02.2022}. В коментариях учащиеся пишут следующее:\textbf{Иванов Иван Б19-021}. После чего, в конце дня, бот должен сделать \textbf{XLS} таблицу и отправить преподавателю, ведущему в этот день занятия. 
 \section{Клиентская часть}
 Здесь я буду перечислять то, что должно быть в клавиатуре клиента (под клавиатурой будем понимать кнопки с командами). Если проекту будет дан зеленый свет, то будет создана схема команд и их зависимости.
  \begin{enumerate}
 \item[•] \textbf{Старт}(start)- приветствие нового пользователя
 \item[•] \textbf{Помощь}(help)- перечисление основных команд
 \item[•] \textbf{Преподаватель сегодня}
 \item[•] \textbf{Спортивные мероприятия}
 \item[•] \textbf{Отключить звук рассылки}
 \item[•] \textbf{Смена языка}(вряд ли сделаю)
 \item[•] \textbf{Количество посещений(отмеченные ботом)}
 \item[•] \textbf{Показать геолокацию занятий}
 \item[•] \textbf{Спортивные новости}
 \item[•] \textbf{Назначить оповещение занятий}
  \end{enumerate}
 + с тектом будут смайлики
  \section{Админская часть}
  В этой части будет функционал, доступный только админам телеграмм-каналов.
    \begin{enumerate}
   \item[•] \textbf{Рассылка сообщения}
    \item[•] \textbf{Создать пост для телеграм-канала}
    \item[•] \textbf{Создание поста о спортивных мероприятиях}
    \item[•] \textbf{Удаления поста о спортивных мероприятиях}
    \item[•] \textbf{Создание поста о спортивных новостях}
    \item[•] \textbf{Удаление поста о спортивных новостях}
    \item[•] \textbf{Создать голосование}
    \item[•] \textbf{Создать объект на карте(геоданные)}
    \item[•] \textbf{Удалить объект на карте(геоданные)}
    \item[•] \textbf{Включить бота)}
    \item[•] \textbf{Отключить бота}
    
	 \end{enumerate}
	  + с тектом будут смайлики
  \section{Работа с другими ботами}
  В телеграмме есть уже созданные беспланые боты для облегчения работы с каналами. Здесь я хочу, чтобы один бот управлял другими.
    \section{Работа с базой данных}
   Стоит заранее продумать о базе данных, чтобы можно было реализовать все клиентские функции. Будет большой проблемой, если бот не ведет учет какого-то нужного параметра клиента( Что-то исправлять будет черезвыйчано сложно ). Также стоит реализовать перевод из \textbf{SQL} в \textbf{XLS}.
     \begin{enumerate}
     \item[•] Создать БЗ для хранения id отправленных и принятых файлов
     	 \end{enumerate}
   \section{Вывод}
	Надеюсь, данный проект будет являться полезным для студентов МФТИ. Возможно, если получится его довести до логического конца, то можно будет сделать подобных ботов и для разных специальностей.   
   Некоторые вещи я уже умею делать, а некоторые даже не представляю как сделать.
\end{document}